\chapter{Discussão}\label{sec:discussao}

Este capítulo interpreta os resultados apresentados no \hyperref[cap4-resultados]{Capítulo 4} à luz da questão de pesquisa formulada na \hyperref[sec-12-problema-de-pesquisa]{Seção 1.2} e das taxonomias de paralelização discutidas na \hyperref[sec-26-gpu-metaheuristicas]{Seção 2.6}.
A discussão está organizada em quatro partes: síntese dos achados e resposta à questão de pesquisa; comparação com o estado da arte; ameaças à validade e limitações do estudo; e implicações práticas para o projeto de algoritmos de roteamento em GPU.

\section{Síntese dos achados}\label{sec-51-sintese-achados}

A questão de pesquisa indaga como comparar, de forma fiel, o impacto de diferentes estratégias de paralelização em GPU sobre o desempenho de um algoritmo memético para o TSP. Os resultados permitem uma resposta em duas partes.
A primeira é que o fator determinante não é o uso da GPU em si, mas \textbf{onde} e \textbf{como} o tráfego de dados entre hospedeiro e dispositivo é organizado: as três variantes aceleradas diferem entre si de forma sistemática, apesar de compartilharem a mesma estrutura algorítmica.
A segunda é que esse efeito é \textbf{multidimensional}: as variantes não se ordenam em uma única escala de mérito, mas trocam posição conforme a métrica considerada.

Sob a ótica do \textbf{tempo de execução}, a variante \textbf{Híbrida Otimizada} é consistentemente superior, com ganho de performance mediano por instância de \(5{,}69\times\) sobre a Híbrida Ingênua e tempo médio agregado de \(18{,}94\) segundos, contra \(47{,}08\) segundos da Ingênua.
Esse resultado é coerente com a hipótese de que o processamento em lote amortiza a latência de transferência: ao enviar a população inteira em uma única chamada, o custo de comunicação deixa de ser proporcional ao número de indivíduos por geração.

\begin{warning}[Resultados provisórios]\label{obs-discussao-provisoria}
  Os números discutidos aqui vêm da execução descrita na \hyperref[obs-ajustes]{Seção 1.3}.
  A limitação específica da variante totalmente em GPU está na \hyperref[obs-fullgpu-iteracoes]{Observação (1)} do \hyperref[cap4-resultados]{Capítulo 4}, e a aritmética dos \emph{gaps}, na \hyperref[obs-aritmetica]{Observação (3)}.
\end{warning}

Sob a ótica da \textbf{qualidade da solução}, a variante \textbf{Totalmente em GPU} ocupou a primeira posição em todas as faixas de tamanho, com \emph{gap} médio agregado de \(0{,}88\%\) contra \(1{,}28\%\) e \(1{,}30\%\) das variantes híbridas.
O teste de Friedman rejeitou a hipótese nula de desempenho equivalente (\(p \approx 6{,}9 \times 10^{-5}\)), e o pós-teste de Nemenyi apontou superioridade da Totalmente em GPU sobre a Ingênua (\(p \approx 0{,}00028\)) e sobre a Otimizada (\(p \approx 0{,}00071\)), sem diferença significativa entre as duas híbridas (\(p \approx 0{,}97\)).

Essa vantagem de qualidade deve ser lida com cuidado por causa da assimetria paramétrica da \hyperref[sec-531-assimetria-parametrica]{Seção 5.3.1}: o orçamento de busca local maior já produz \emph{gap} menor por si só.
Como a variante Híbrida Otimizada também evita transferências repetidas, via processamento em lote, e mesmo assim não alcança a mesma qualidade, o número de iterações parece ser o fator determinante, e não a permanência dos dados na GPU.

Um segundo eixo de análise, exigido pelo enquadramento do problema, é o do \textbf{comportamento de memória}, que é distinto do tempo de execução e não se reduz a ele. A Tabela~\ref{tbl:memory-axis} sintetiza o perfil de transferência e a pegada de memória de cada variante.
Como a matriz de distâncias é armazenada completa e em inteiros de 32 bits, ocupando \(n_{coords}^2 \times 4\) bytes, e a população tem tamanho fixo de \(256\) indivíduos, o termo quadrático da matriz domina o termo linear da população já em instâncias pequenas.
O fator que limita o porte das instâncias tratáveis em um dispositivo de \(4\) GB é, portanto, o armazenamento quadrático das distâncias (\hyperref[sec-311-hardware-software]{Seção 3.1.1}), e não o operador de busca local.
O ponto central é que velocidade e capacidade disputam o mesmo recurso físico: a estratégia que minimiza o tempo de comunicação é justamente a que concentra o estado na memória do dispositivo.

Cabe registrar que esse gargalo de memória não é uma barreira sem soluções na literatura moderna, e as principais alternativas são conhecidas.
A mais direta é o \textbf{armazenamento triangular compactado}, que explora a simetria do TSP e reduz a matriz de \(n(n-1)/2\) entradas, exatamente metade do custo.
Uma segunda família de abordagens troca memória por computação, \textbf{recalculando distâncias sob demanda} em vez de materializá-las, o que troca o consumo quadrático por um fator constante de trabalho aritmético.
Uma terceira linha trabalha com \textbf{particionamento da vizinhança} \emph{(neighbourhood splitting)}: Schulz e colaboradores descrevem explicitamente o uso de memória como o critério que determina o tamanho do subconjunto de movimentos avaliado por vez, relatando limiares concretos em que a busca \emph{k}-opt deixa de caber e precisa ser fatiada \autocite{schulz2013gpu}.
Por fim, estratégias de \textbf{transferência em blocos} permitem processar instâncias cujo estado completo não caberia simultaneamente na memória do dispositivo.

\begin{longtable}[]{@{}
    >{\raggedright\arraybackslash}p{(\linewidth - 4\tabcolsep) * \real{0.2222}}
    >{\raggedright\arraybackslash}p{(\linewidth - 4\tabcolsep) * \real{0.3889}}
    >{\raggedright\arraybackslash}p{(\linewidth - 4\tabcolsep) * \real{0.3889}}@{}}
  \caption{Eixo de memória e comunicação por variante (análise estrutural,
    não medida experimental)
    \label{tbl:memory-axis}}\tabularnewline
  \toprule\noalign{}
  Variante                   & Perfil de transferência por geração
                             & Estado residente no dispositivo                                                                            \\
  \midrule\noalign{}
  \endfirsthead
  \toprule\noalign{}
  Variante                   & Perfil de transferência por geração
                             & Estado residente no dispositivo                                                                            \\
  \midrule\noalign{}
  \endhead
  \bottomrule\noalign{}
  \endlastfoot
  \textbf{CPU}               & nenhuma                                                       & nenhum                                     \\
  \textbf{Híbrida-Ingênua}   & \(O(n_{coords})\) por indivíduo, em ambas as direções         & apenas o indivíduo em execução             \\
  \textbf{Híbrida-Otimizada} & \(O(n_{pop} \cdot n_{coords})\) por geração, agrupada em lote & o lote avaliado                            \\
  \textbf{Totalmente em GPU} & \(O(n_{coords}^2)\) uma única vez, na inicialização           & população e matriz de distâncias completas \\
\end{longtable}

A ordem de grandeza da penalidade de comunicação é dada pelo contraste entre a largura de banda do barramento PCIe (\(\approx 16\) GB/s) e a da memória interna do dispositivo (\(\approx 112\) GB/s) (\hyperref[sec-311-hardware-software]{Seção 3.1.1}): uma transferência que atravesse o barramento custa cerca de sete vezes o mesmo volume movido internamente.
Essa razão explica por que a variante Ingênua, mesmo executando a busca local em GPU, não se distancia das demais variantes em tempo, e por que a Otimizada, ao reduzir o número de travessias, converte a capacidade de cálculo do dispositivo em ganho efetivo.

\section{Comparação com o estado da arte}\label{sec-52-comparacao-estado-arte}

Embora as abordagens paralelas desenvolvidas tenham alcançado alta velocidade de execução, é importante situar as limitações do GA+\(2\)-opt implementado frente ao estado da arte contemporâneo no TSP.
Algoritmos dedicados e altamente especializados, notadamente o algoritmo LKH-3 \autocite{helsgaun2000effective} e o algoritmo genético baseado em Edge Assembly Crossover (EAX) \autocite{nagata2013powerful}, rotineiramente atingem \(0\%\) de \emph{gap} (garantindo a solução ótima) na esmagadora maioria das instâncias padrão da TSPLIB.

O GA+\(2\)-opt aqui explorado adota operadores canônicos e uma arquitetura base focada na investigação didática das relações de custo-benefício da programação paralela.
Consequentemente, este trabalho foca na comparação entre as variantes do mesmo algoritmo (analisando os ganhos estruturais de cada estratégia de implantação), e não na competição isolada por poder de refinamento contra as meta-heurísticas de última geração que utilizam buscas como \(5\)-opt ou estratégias dinâmicas complexas de vizinhança.
Uma investigação futura poderia avaliar a integração de operadores mais sofisticados (como o LK-move) às implementações em GPU, visando estreitar essa lacuna de qualidade.

Essa delimitação é relevante para a leitura correta dos resultados: o \emph{gap} de \(0{,}88\%\) a \(1{,}30\%\) reportado no \hyperref[cap4-resultados]{Capítulo 4} não deve ser lido como uma medida de competitividade do estado da arte, mas como um indicador \emph{relativo} da capacidade de refinamento de cada estratégia de paralelização dado o mesmo núcleo algorítmico.
A contribuição do trabalho está na comparação controlada entre variantes, não no desempenho absoluto frente a resolvedores especializados.

\section{Ameaças à validade e limitações}\label{sec-53-ameacas-validade}

\subsection{Assimetria paramétrica entre variantes}\label{sec-531-assimetria-parametrica}

A limitação mais severa deste estudo é a divergência no limite de iterações da busca local \(2\)-opt entre as variantes: a versão totalmente em GPU foi executada nos experimentos originais com \(50\) varreduras por indivíduo, enquanto as variantes em CPU e híbridas utilizaram \(10\) (\hyperref[sec-321-ga-base]{Seção 3.2.1}).
Essa assimetria beneficia diretamente a variante que obteve o melhor \emph{gap} e, portanto, compromete a interpretação desse resultado como efeito da residência dos dados na GPU.

\subsection{Assimetria de semântica entre os operadores \texorpdfstring{\(2\)-opt}{2-opt}}\label{sec-5311-semantica-operadores}

Além do limite de iterações, há uma diferença que a normalização do parâmetro \textbf{não} corrige, porque está na definição dos operadores.
Os quatro \emph{kernels} de busca local não fazem o mesmo trabalho: a variante \textbf{Totalmente em GPU} varre a vizinhança até não haver melhoria, aceitando a primeira melhoria encontrada (\emph{first-improvement}); a implementação em \textbf{CPU} avalia a varredura completa e aplica a melhor melhoria (\emph{best-improvement}); e as variantes \textbf{híbridas} executam uma única passada por chamada de \emph{kernel}.
Ou seja, um mesmo valor de \(n_{2\text{-opt}}\) representa quantidades diferentes de avaliações e trajetórias de descida distintas.

% TODO Ask about this later. Didn't understand exactly what's happening.

Essa diferença é mais difícil de contornar do que a anterior, porque decorre do modo como cada operador foi implementado, e não de um parâmetro mal propagado.
Para o desenho experimental deste trabalho, a consequência prática é que as variantes compartilham a estrutura algorítmica, e a comparação entre elas mede também a diferença entre os operadores — não apenas o efeito de executar na CPU ou na GPU.

Um efeito correlato, também de natureza estrutural, é que a variante Totalmente em GPU contabiliza um número de gerações bem maior que a Híbrida Otimizada em instâncias do mesmo porte, o que reflete diferenças na forma de contar gerações no dispositivo e no hospedeiro.
Esse fator concorre com os anteriores na explicação das diferenças de tempo e de qualidade observadas.

\subsection{Instrumentação de desempenho ausente}\label{sec-532-instrumentacao}

A análise baseia-se exclusivamente em medições macroscópicas de tempo total de execução, sem instrumentação por \emph{kernel}.
Ferramentas de \emph{profiling} de hardware como o \emph{NVIDIA Nsight Compute}, embora capazes de fornecer contadores detalhados (ocupação, pressão de registradores, divergência de \emph{warps}), serializam a execução dos kernels e exigem múltiplos passes de repetição, distorcendo o tempo real em algoritmos iterativos de longa duração como o GA aqui implementado \autocite{nvidia2024cuda}.
Como consequência, as explicações arquiteturais oferecidas no \hyperref[cap4-resultados]{Capítulo 4} para a degradação de desempenho em instâncias grandes — pressão de registradores e divergência de \emph{warps} — permanecem como hipóteses plausíveis, porém não verificadas. Uma abordagem menos intrusiva, baseada em temporizadores via primitivas de eventos CUDA na camada Python/CuPy, é o caminho natural para instrumentação futura.

\subsection{Hardware único e validade externa}\label{sec-533-hardware-unico}

Todos os experimentos foram conduzidos em um único dispositivo, uma \textbf{NVIDIA GeForce GTX 1050 Mobile} com \(4\) GB de memória (\hyperref[sec-311-hardware-software]{Seção 3.1.1}).
Trata-se de uma arquitetura de entrada, com número reduzido de multiprocessadores e largura de banda de memória limitada.
Embora essa escolha tenha sido metodologicamente vantajosa por tornar as ineficiências de implementação mais evidentes do que seriam em aceleradores de alto desempenho, ela restringe a generalidade das conclusões: os limiares de \(n_{coords}\) em que a estratégia em lote deixa de compensar as transferências, e o ponto em que a variante residente passa a sofrer degradação, são específicos deste dispositivo e não devem ser extrapolados para GPUs de data center sem verificação empírica.

\subsection{Linha de base incompleta}\label{sec-534-baseline-incompleta}

A variante puramente em CPU foi executada apenas para instâncias com \(n_{coords} \le 100\), devido ao tempo proibitivo de execução nas demais (\hyperref[sec-33-selecao-instancias]{Seção 3.3}).
Por consequência, para instâncias médias e grandes não existe referência em CPU no conjunto de dados, e as comparações de ganho de performance tomam a Híbrida-Ingênua como base. Isso impede quantificar o ganho absoluto da aceleração em GPU para \(n_{coords} > 100\) e desloca o sentido das comparações para o interior do conjunto de variantes aceleradas.

\subsection{Métricas e critérios de parada}\label{sec-535-metricas-criterios}

As métricas adotadas são limitadas a tempo total de execução e qualidade da solução final.
Não foram medidos consumo energético, uso de memória ao longo da execução nem tempo por fase do algoritmo, de modo que não é possível atribuir as diferenças observadas a componentes específicos do ciclo evolutivo. Além disso, os critérios de parada dependem de parâmetros de paciência e de limite de gerações proporcionais a \(n_{coords}\) (\hyperref[sec-321-ga-base]{Seção 3.2.1}), o que faz com que instâncias maiores recebam orçamento computacional proporcionalmente maior; comparações entre faixas de tamanho devem ser lidas com essa ressalva.
Por fim, a equivalência entre variantes é \emph{estrutural}, e não bit a bit: diferenças na associatividade de reduções em ponto flutuante e condições de contorno de hardware impedem determinismo perfeito entre CPU e GPU (\hyperref[sec-32-arquitetura-variantes]{Seção 3.2}).

\section{Implicações práticas}\label{sec-54-implicacoes-praticas}

Do ponto de vista de quem projeta ou emprega algoritmos de roteamento em GPU, os resultados sugerem uma regra de decisão orientada por aplicação, e não por uma hierarquia única de variantes:

\begin{itemize}
  \tightlist
  \item
        \textbf{Quando a latência domina} — cenários interativos, planejamento com restrição de tempo de resposta ou execução de muitas instâncias em sequência — a \textbf{Híbrida Otimizada} é a escolha indicada.
        Ela oferece o menor tempo de execução observado, com ganhos de até \(616{,}9\times\) sobre a CPU em instâncias pequenas, ao custo de manter o estado principal no hospedeiro e pagar uma transferência em lote por geração.
  \item
        \textbf{Quando a qualidade da rota domina} — e o tempo de resposta, embora relevante, é secundário — a variante \textbf{Totalmente em GPU} é a candidata natural, por concentrar o estado no dispositivo e evitar travessias do barramento durante a evolução.
  \item
        \textbf{Quando a instância cresce} — aproximando-se do teto de memória do dispositivo — a variante residente perde vantagem, e a estratégia em lote passa a ser a única que preserva simultaneamente tempo competitivo e uso de memória controlado.
        Nesse regime, a decisão deixa de ser sobre velocidade e passa a ser sobre viabilidade.
\end{itemize}

Um corolário prático é que a escolha da estratégia de paralelização deve ser tratada como parte do projeto do algoritmo, e não como um detalhe de implementação.
As variantes compartilham a mesma estrutura algorítmica e diferem sobretudo na organização da execução e do tráfego de dados, o que mostra que essas decisões de arquitetura merecem atenção explícita em estudos comparativos de meta-heurísticas em GPU.
Vale lembrar que essa atribuição é aproximada: como a \hyperref[sec-5311-semantica-operadores]{Seção 5.3.1} descreve, os operadores de busca local não são idênticos entre as variantes, de modo que parte das diferenças observadas vem também deles.


